\section{Related Work}
\label{sec:related_work}

\begin{table*}[t]
\centering
\caption{Comparison of DANTE with representative baselines across the key design dimensions of decentralized personalized learning in vehicular networks.}
\label{tab:comparison}
\renewcommand{\arraystretch}{1.15}
\small
\begin{tabular}{l c c c c c c}
\hline
\textbf{Method}
  & \textbf{Adaptive}
  & \textbf{Personalized}
  & \textbf{Resource}
  & \textbf{Dual}
  & \textbf{Byzantine}
  & \textbf{Learned} \\
  & \textbf{Topology}
  &
  & \textbf{Aware}
  & \textbf{Interface}
  & \textbf{Robust}
  & \textbf{Selection} \\
\hline
FedAvg~\cite{mcmahan2017communication}
  & \texttimes & \texttimes & \texttimes
  & \texttimes & \texttimes & \texttimes \\
D-PSGD~\cite{lian2017dsgd}
  & \texttimes & \texttimes & \texttimes
  & \texttimes & \texttimes & \texttimes \\
CHOCO-SGD~\cite{koloskova2019choco}
  & \texttimes & \texttimes & \checkmark
  & \texttimes & \texttimes & \texttimes \\
BALANCE~\cite{fang2024balance}
  & \texttimes & \texttimes & \texttimes
  & \texttimes & \checkmark & \texttimes \\
MDFL~\cite{chen2025mobility}
  & \checkmark & \texttimes & \checkmark
  & \texttimes & \texttimes & \checkmark \\
CDPL~\cite{boubouh2020robust}
  & \checkmark & \checkmark & \texttimes
  & \texttimes & \checkmark & \checkmark \\
DPFL~\cite{kharrat2025dpfl}
  & \checkmark & \checkmark & \checkmark
  & \texttimes & \texttimes & \texttimes \\
pFedGraph~\cite{ye2023pfedgraph}
  & \checkmark & \checkmark & \texttimes
  & \texttimes & \texttimes & \texttimes \\

\hline
\textbf{DANTE (ours)}
  & \checkmark & \checkmark & \checkmark
  & \checkmark & \checkmark & \checkmark \\
\hline
\end{tabular}
\end{table*}


The growing demand for privacy-preserving, low-latency collaborative intelligence at the network edge has made decentralized learning a major research direction, with a large body of work now addressing serverless optimization, personalized collaboration, communication efficiency, and topology design. Early decentralized federated optimization focused on replacing server-based aggregation with peer-to-peer mixing. D-PSGD~\cite{lian2017dsgd} established the canonical gossip-based training paradigm in which clients alternate local SGD and neighbor averaging over a fixed mixing graph. D$^2$~\cite{tang2018d2} strengthened this line by removing the dominant penalty induced by data heterogeneity, while SGP~\cite{assran2019sgp} extended decentralized training to directed and time-varying communication graphs through PushSum-based gossip. CHOCO-SGD~\cite{koloskova2019choco} further reduced communication overhead by combining decentralized optimization with compressed message exchange. These methods provide the algorithmic foundation of decentralized federated learning, but they still optimize a shared objective over generic peer graphs and therefore do not address personalized collaboration.

This limitation becomes particularly restrictive in vehicular networks, where clients are mobile, communication opportunities are short-lived, and local data distributions are strongly context dependent. Barbieri et al.~\cite{barbieri2022extended} demonstrated that consensus-driven decentralized federated learning can support extended sensing in connected vehicles through direct V2V model exchange. Yuan et al.~\cite{yuan2023p2p} proposed peer-to-peer federated continual learning for naturalistic driving action recognition, showing that decentralized exchange can preserve privacy while mitigating catastrophic forgetting. Chen et al.~\cite{chen2025mobility} moved closer to practical V2X deployment by optimizing local iteration counts and leader selection under mobility and wireless constraints through multi-agent reinforcement learning. Although these studies bring decentralized learning into vehicular settings, they primarily target efficient coordination of shared-model training and do not jointly address personalization, dual link mode for efficient and low latency communication, and Byzantine-aware collaboration.

Because vehicular data are inherently non-IID, personalized decentralized learning is a more natural formulation than consensus on a single global model. Various works~\cite{vanhaesebrouck2017decentralized,boubouh2020robust,basmadjian2022advantages} explored the concept of personalized learning over peer-to-peer networks. \cite{boubouh2020robust} proposed CDPL, a robust and adaptive algorithm to learn personalized models in dynamic networks by evaluating the received updates against a small subset of the peer's local dataset. The authors in~\cite{zantedeschi2020fully} suggested to jointly learn personalized models and sparse collaboration graphs in a fully decentralized setting. DisPFL~\cite{dai2022dispfl} introduced decentralized sparse training and personalized masks, reducing both communication and computation while maintaining client-specific models. DPFL~\cite{kharrat2025dpfl}, which we also implement, cast decentralized personalization as a bilevel optimization problem and used constrained greedy graph construction to select collaborators under a fixed budget. These methods establish that who collaborates with whom is central to personalization quality, yet they do not account for the rapidly changing neighborhoods, asymmetric interface costs, and trust uncertainty that define vehicular V2X environments.

Operational constraints add a second layer of difficulty. Communication and computation efficiency have been studied through compression, resource allocation, and graph sparsification. CHOCO-SGD~\cite{koloskova2019choco} reduces payload size through compressed gossip, while Meng et al.~\cite{meng2025resource} jointly optimize neighbor selection, transmit power, and CPU frequency using graph neural networks in decentralized wireless learning. Zhang et al.~\cite{zhang2024energy} showed that graph sparsification can substantially reduce energy consumption at bottleneck nodes while preserving learning quality. Closely related to these efforts is the topology-optimization literature, which treats the collaboration graph itself as a statistical design variable. pFedGraph~\cite{ye2023pfedgraph} infers collaboration weights from model similarity and dataset size to produce client-specific aggregation patterns. DFedPGP~\cite{liu2024dfedpgp} showed that directed collaboration can outperform symmetric peer exchange in personalized decentralized learning. Le Bars et al.~\cite{lebars2023topology} further demonstrated that topology should be adapted to data heterogeneity itself, and formalized topology learning as an optimization problem driven by neighborhood heterogeneity. Together, these works show that resource allocation and topology design are decisive for decentralized learning performance, but they stop short of unifying personalization, hybrid V2X link heterogeneity, and adversarial robustness in a single decentralized control framework.

DANTE differs from all of these approaches in three coupled ways. First, it formulates vehicular collaboration as a dynamic network formation problem rather than fixed gossip, centralized graph inference, or static budgeted graph construction. Second, it explicitly differentiates the cost and reachability of PC5 sidelink and Uu relay communication, allowing each vehicle to trade off local low-cost exchange against longer-range but more expensive collaboration. Third, it embeds trust estimation and robust aggregation directly into topology control, so collaborator selection depends jointly on personalized learning value, communication cost, and reliability. To the best of our knowledge, no prior method combines these properties in a fully decentralized personalized learning framework for hybrid 5G C-V2X networks.


